\section{Expected performance and limits of validity}
\label{sec:expectations}

\subsection{What to expect}

\begin{itemize}
\item \textbf{Bright stars are scintillation-limited.} For the 358-mm
  reference telescope, a $V\simeq5$ star in \SI{60}{s} reaches
  $\mathrm{S/N}\simeq1.3\times10^{3}$, not the $\sim2\times10^{4}$ a
  photon-only budget would claim (Fig.~\ref{fig:snrbudget}); the ceiling
  scales as $d^{2/3} X^{-1.75} \sqrt{t}$.
\item \textbf{Faint limits are sky-limited.} Beyond
  $m \gtrsim \mu_{\rm sky} - 2.5\log_{10}(\pi r^2)$ the S/N falls as
  $10^{-0.4m}$ and doubling the exposure gains only $\sqrt{2}$.
\item \textbf{Moonlight matters within $\sim90^\circ$.} Full Moon at
  $30^\circ$ separation brightens the $V$ sky by $\simeq4$--5 mag
  (Sect.~\ref{sec:sky}), turning background-limited exposures from minutes
  into hours.
\item \textbf{Slit spectroscopy away from the parallactic angle} loses tens
  of per cent at the blue end already at $X = 1.5$--2
  (Fig.~\ref{fig:dispersion}); the Results window reports $q$ at every
  sample so the slit can be set correctly.
\item \textbf{Star Analyser work} is photometrically generous (no slit
  losses) but limited to $R\sim100$--200 by seeing; S/N per element of
  several hundred in minutes is realistic for $V \lesssim 8$
  (Fig.~\ref{fig:sa100}).
\end{itemize}

\subsection{Known approximations}

The model is explicit about its accuracy floor: the $U$--$K$ sky table and
its colour interpolation are broad-band, and the twilight/daylight bridge is
approximate; the extinction fallback table is a smooth broad-band baseline
until a site-specific curve is loaded; scintillation uses the classical
Young coefficient (site-specific coefficients from \citealt{osborn2015} can
be substituted); atmospheric dispersion in slit mode assumes the worst-case
orientation; extended sources are uniform discs; and slitless mode awaits
validation against a calibrated instrument. Flat-field and fringing
residuals, cosmic rays and sky-subtraction penalties are not modelled.

\subsection{Numerical verification}

Two test programs ship with the code: a synthetic smoke test that pins the
scintillation-limited photometric S/N and the per-element spectroscopic S/N
of the reference configuration, and a regression suite that verifies, among
others, Eq.~(\ref{eq:ks}) against the published scattering function, the
ecliptic symmetry of the zodiacal term, Eq.~(\ref{eq:pickering}) against
$\sec z$, the Student-$t$ Moffat coupling, the Filippenko dispersion scale,
Eq.~(\ref{eq:sa}) for the SA100/SA200, and the extended-source area
fractions.
